% Synchronized proof export for TEXTBOOK-PART-IV-CHAPTER-14-INFORMATION-THEORY-SPINE
\section{Chapter 14 Information-Theory Spine}

\paragraph{Status.}
This note is a synchronized mathematical export for task
\texttt{TEXTBOOK-PART-IV-CHAPTER-14-INFORMATION-THEORY-SPINE}.  The scoped
Section~14.2 declaration surface compiles at Lean source commit
\texttt{10dfd88}.  It covers Eqs.~(14.4)--(14.6), Theorem~14.1, the event
specialization of Exercise~14.10, and Theorem~14.2.  It does not claim the
entropy/coding results of Section~14.1, full sub-sigma-algebra data processing,
or an adaptive bandit-history chain rule.

\paragraph{Lean declarations.}
The public surface is
\texttt{relativeEntropy}, its regular, probability, singular, and finiteness
adapters; \texttt{bernoulliRelativeEntropy}; the RN-restriction and KL-split
helpers; \texttt{bernoulliKLCore\_event\_le}; the binary affinity and overlap
lemmas; \texttt{binaryBretagnolleHuberCore};
\texttt{bretagnolleHuberScale} and its nonnegativity;
\texttt{binaryBretagnolleHuber};
\texttt{bernoulliRelativeEntropy\_event\_le}; scale antitonicity; and the
terminal \texttt{bretagnolleHuber}.

\paragraph{Measure and event reduction.}
Define $D(P\Vert Q)$ as Mathlib's extended-real KL divergence.  On the regular
branch $P\ll Q$ with an integrable log likelihood ratio, it is the nonnegative
extended-real embedding of the RN/log-likelihood integral; the finite-measure
mass correction vanishes for probability laws.  Failure of absolute
continuity gives infinity, and finiteness is equivalent to absolute continuity
plus integrability.

For a measurable event $A$, restrict both laws to $A$ and $A^{\mathrm c}$.
The restricted RN derivative agrees almost everywhere with the original one,
so the KL integral splits exactly over the two cells.  Applying the convex
mass-level \texttt{klFun} bound to both cells gives
\[
  d(P(A),Q(A))\leq D(P\Vert Q).
\]
The endpoint cases $Q(A)\in\{0,1\}$ use absolute continuity, and the
infinite-KL case is immediate.  The direction is always $P$ to $Q$.

\paragraph{Binary testing.}
For interior $q$, concavity of logarithm gives
\[
  e^{-d(p,q)/2}
  \leq \sqrt{pq}+\sqrt{(1-p)(1-q)}.
\]
A two-term square-root inequality bounds one half of the squared affinity by
$p+(1-q)$.  After squaring the exponential inequality,
\[
  \frac12e^{-d(p,q)}\leq p+(1-q).
\]
The cases $q=0,1$ are discharged through the exact Bernoulli support
convention.  The real scale is defined as $e^{-d}/2$ for finite extended-real
$d$ and zero at $d=\infty$; it is nonnegative and antitone.

\paragraph{Source terminal and boundary.}
Event data processing, scale antitonicity, the endpoint-complete binary
theorem, and $Q(A^{\mathrm c})=1-Q(A)$ yield
\[
  \operatorname{BHScale}(D(P\Vert Q))
  \leq P(A)+Q(A^{\mathrm c}).
\]
This is the unconditional Theorem~14.2 terminal.  It has no finite-KL premise,
mutual absolute-continuity premise, policy, horizon, filtration, kernel chain
rule, or stopping-time assumption.  Chapter~15 must still build the
same-policy adaptive-history laws and divergence decomposition that consume
this theorem.
